\paragraph{The primitive exterior-square bridge used in the proof.}
Here is an elementary construction of the classical five-dimensional
orthogonal representation, with the precise unipotent comparison.
Let $E=\mathbb C^4$ have a symplectic basis $e_1,e_2,f_1,f_2$.
The invariant bivector $\omega=e_1\wedge f_1+e_2\wedge f_2$ has
nonzero square. The wedge pairing on $\bigwedge^2E$ is symmetric
and nondegenerate. Its orthogonal complement
$E_0=\omega^\perp$ therefore has dimension five and is nondegenerate.
Preservation of $\omega^2$ gives determinant one on $E$;
$\det(\bigwedge^2g)=(\det g)^3$ and the fixed line
$\mathbb C\omega$ therefore give determinant one on $E_0$.
The action gives
\[
 \rho:\operatorname{Sp}(E)\longrightarrow\operatorname{SO}(E_0).
\]
Its differential is injective. Indeed, a symplectic infinitesimal
transformation acting trivially on $E_0$ also acts trivially on
$\mathbb C\omega$, hence on all of $\bigwedge^2E$. In a basis,
the equations $Xe_i\wedge e_j+e_i\wedge Xe_j=0$ for every $i<j$
first make every off-diagonal coefficient of $X$ zero, and then give
$X_{ii}+X_{jj}=0$ for every distinct pair. Since there are four
indices, these equations make every diagonal coefficient zero.
Both Lie algebras have dimension ten, so the image contains a
neighbourhood of the identity in $\operatorname{SO}(E_0)$.
That complex special orthogonal group is connected. For completeness,
orthogonal elimination expresses an orthogonal matrix as a product
of reflections in vectors of nonzero norm: to fix a nonisotropic
vector, reflection in the difference of its image and itself works
when that difference is nonisotropic; otherwise reflection in their
sum followed by reflection in the vector works. At least one of
sum and difference is nonisotropic. Continue on its nondegenerate
orthogonal complement. Determinant one makes the number of
reflections even. The nonisotropic vectors of a complex vector
space of dimension at least two form a path-connected set, since
two generic complex line segments can avoid the zero set of the
norm polynomial except at finitely many parameter values, around
which the paths can be detoured. Varying one vector in a pair of
reflections to the other connects that pair to the identity.
The one-dimensional final complement contributes only the already
accounted-for sign. This proves connectedness. Thus the open image
of $\rho$ is the whole special orthogonal group.

The kernel is $\{I,-I\}$. If an element acts trivially on
$E_0$, it also fixes $\omega$ and hence every bivector.
It preserves every two-plane, as follows by applying it to a
nonzero decomposable bivector of that plane. Intersections of
two-planes show that it preserves every line. A linear map
preserving every line is scalar: apply it to the coordinate lines
and their pairwise sums. Its scalar has square one. In particular
every finite orthogonal element of order $m$ has lifts whose
$m$th powers belong to $\{I,-I\}$; the lifts are finite and
semisimple.

The symplectic transvection fixing $e_1,e_2,f_2$ and sending
$f_1$ to $f_1+e_1$ acts on the following primitive basis:
\[
 e_1\wedge e_2,\quad e_1\wedge f_2,\quad
 e_2\wedge f_1,\quad f_1\wedge f_2,\quad
 e_1\wedge f_1-e_2\wedge f_2.
\]
Its difference from the identity sends the third vector to
the negative of the first and the fourth to the second, and
annihilates the other three. Its orthogonal Jordan type is
therefore $(2,2,1)$. Conversely, for an orthogonal unipotent
$I+N$ of that type, $N^2=0$ implies $N^*=-N$, because
$(I+N)^*=(I+N)^{-1}=I-N$. Its image is an isotropic plane.
After choosing a basis of that plane and a dual isotropic
plane, use the isometry on their residual nondegenerate line
to make the identification of two such decompositions have
determinant one. Changing its sign changes that determinant
and commutes with $N$, which kills the residual line.
In these decompositions, $N$ is determined by a nonzero alternating form on
the dual two-dimensional space. All such forms are equivalent
under its general linear group. Extending that change by the
inverse transpose on the first plane gives an orthogonal
change of basis; its determinant is one. Thus these unipotents
form one special-orthogonal conjugacy class. Surjectivity of
$\rho$ lifts the conjugacy. The two lifts of $I+N$ are a
symplectic transvection and its negative. Hence, after changing
one lift's sign, the product of the two finite lifts in the
argument is precisely a rank-one symplectic transvection.

Finally a symplectic diagonal element with eigenvalues
$x,x^{-1},y,y^{-1}$ acts on $E_0$ with eigenvalues
\[
 1,\quad xy,\quad x/y,\quad y/x,\quad (xy)^{-1}.
\]
This gives the exact eigenvalue comparison used above and
completes the required bridge without changing any assertion
about rational or integral conjugacy.
